%% ============================================================================
%% theory_fragments.tex  --  fragments for the round-2 revision (Reviewer #1,
%% points 1, 3, 5, 6).  amsart style, notation and labels of FBS-rev2-arxiv.tex.
%%
%% Needed in the preamble (the paper only defines theorem/definition/remark):
%%   \theoremstyle{plain}
%%   \newtheorem{proposition}{Proposition}
%%   \newtheorem{lemma}{Lemma}
%%   \newtheorem{corollary}{Corollary}
%%   \DeclareMathOperator{\dist}{dist}     % the paper uses \operatorname{dist}
%% Placeholders to be filled by other agents: \texttt{[RUNTIME-TBD]},
%% \texttt{[MU-TABLE-TBD]}, \texttt{[PLAP-TBD]}.
%% New bibliography entries are collected at the end of this file.
%% The fragments are independent; each block says where it is meant to go.
%% ============================================================================


%% ============================================================================
%% FRAGMENT 0 -- standing notation (put at the beginning of the new
%% "Analysis of the method" section, or merge into Section 2)
%% ============================================================================

\subsection*{Setting and notation}

Throughout this section $L=-\Delta+V$ denotes the self-adjoint realization of the
operator in \eqref{eigen} on $H=L^2(\Omega)$ with homogeneous Dirichlet boundary
conditions, with domain $D(L)$; we assume that $L$ is bounded from below and has compact
resolvent (this is the case for $\Omega$ bounded and $V\in L^\infty(\Omega)$, where
$D(L)=H^2(\Omega)\cap H^1_0(\Omega)$ if $\Omega$ is convex or of class $C^{1,1}$).
We denote by $\langle\cdot,\cdot\rangle$ and $\|\cdot\|$ the inner product and norm of
$L^2(\Omega)$, by $E_1\le E_2\le\cdots$ the eigenvalues of $L$ repeated according to
multiplicity, by $\{u_k\}_{k\in\N}$ an orthonormal basis of eigenfunctions, by
$\sigma=\{E_k\}_{k\in\N}$ the spectrum, and
\[
d(E):=\dist(E,\sigma)=\min_k|E-E_k| .
\]
With this notation \eqref{Lambda} reads $\Lambda(u,E)=\|(L-E)u\|^2/\|u\|^2$ and
\eqref{cota.igual} reads $\Lambda(E)=d(E)^2$.
All the statements below remain valid verbatim if the Lebesgue measure is replaced by the
normalized measure $dx/|\Omega|$, which is what a Monte Carlo average with uniformly
distributed points approximates.


%% ============================================================================
%% FRAGMENT 1 -- answers R1.3 (choice of mu) and part of R1.6 (spurious minima)
%% Suggested place: right after the definition of the penalized loss (lossmu),
%% replacing the heuristic paragraph "One technical point ... scale mu ~ E^2".
%% ============================================================================

\subsection*{The penalized loss: exact loss curve and absence of spurious minima}

In practice we do not minimize the quotient \eqref{Lambda} but the penalized loss
\eqref{lossmu}, that is,
\[
\bar\Lambda_\mu(u,E)=\|(L-E)u\|^2+\mu\big(\|u\|^2-1\big)^2,\qquad u\in D(L).
\]
Its infimum, its minimizers and all its critical points can be computed explicitly.

\begin{proposition}\label{prop.penalized}
Let $\mu>0$, $E\in\R$ and $d=d(E)$. Let $N_E$ be the (finite dimensional) linear span of
the eigenfunctions $u_k$ with $|E_k-E|=d$.
\begin{enumerate}
\item[(a)] (Loss curve.)
\[
\bar\Lambda_\mu(E):=\inf_{u\in D(L)}\bar\Lambda_\mu(u,E)=
\begin{cases}
d^2-\dfrac{d^4}{4\mu} & \text{if } d^2\le 2\mu,\\[2mm]
\mu & \text{if } d^2\ge 2\mu .
\end{cases}
\]
\item[(b)] (Minimizers.) If $d^2<2\mu$, the minimizers are exactly the functions
$u=\sqrt{s_*}\,w$ with $w\in N_E$, $\|w\|=1$ and
\[
s_*=\|u\|^2=1-\frac{d^2}{2\mu}\in(0,1].
\]
For them the two terms of the loss are
$\|(L-E)u\|^2=d^2\big(1-\frac{d^2}{2\mu}\big)$ and
$\mu(\|u\|^2-1)^2=\frac{d^4}{4\mu}$, and $\Lambda(u,E)=d^2$.
If $d^2\ge 2\mu$, the unique minimizer is the trivial function $u=0$.
\item[(c)] (Critical points.) The critical points of $u\mapsto\bar\Lambda_\mu(u,E)$ in
$D(L)$ are $u=0$ and the functions $u=\sqrt{1-\lambda/(2\mu)}\,w$, where
$\lambda\in\{(E_k-E)^2\colon k\in\N\}$, $\lambda<2\mu$, and $w$ is a unit vector in the
span of the eigenfunctions $u_k$ with $(E_k-E)^2=\lambda$. Among them, the only local
minimizers are the global minimizers described in (b). In particular, $u=0$ is a local
minimizer if and only if $d^2\ge 2\mu$.
\end{enumerate}
\end{proposition}

\begin{proof}
Write $u=\sum_k c_ku_k$ and $s=\|u\|^2=\sum_kc_k^2$. Then
\begin{equation}\label{expansion}
\|(L-E)u\|^2=\sum_k (E_k-E)^2c_k^2\ \ge\ d^2 s ,
\end{equation}
with equality if and only if $u\in N_E$. Hence
$\bar\Lambda_\mu(u,E)\ge g(s):=d^2s+\mu(s-1)^2$, with equality if and only if $u\in N_E$.
Since $g'(s)=d^2-2\mu+2\mu s$, if $d^2<2\mu$ the function $g$ attains its minimum over
$[0,\infty)$ only at $s_*=1-d^2/(2\mu)>0$ and
$g(s_*)=d^2s_*+\mu\,d^4/(4\mu^2)=d^2-d^4/(4\mu)$, while if $d^2\ge2\mu$ then $g$ is
strictly increasing on $[0,\infty)$ and its minimum is $g(0)=\mu$. This proves (a) and (b)
(for $d^2=2\mu$ both expressions in (a) agree).

(c) For $u,v\in D(L)$,
\[
\frac{d}{dt}\bar\Lambda_\mu(u+tv,E)\Big|_{t=0}
=2\langle (L-E)u,(L-E)v\rangle+4\mu(s-1)\langle u,v\rangle .
\]
Taking $v=u_k$ we see that $u$ is a critical point only if
$\big[(E_k-E)^2-2\mu(1-s)\big]c_k=0$ for every $k$. Thus either $u=0$, or all the indices
$k$ with $c_k\neq0$ share the same value $\lambda=(E_k-E)^2$, and
$s=1-\lambda/(2\mu)$, which forces $\lambda<2\mu$. Conversely, such a $u$ is a finite
combination of eigenfunctions with $(L-E)^2u=\lambda u$, and the derivative above equals
$\big(2\lambda+4\mu(s-1)\big)\langle u,v\rangle=0$ for every $v$.

Let now $u\ne0$ be a critical point with $\lambda>d^2$ and let $u_j$ be an eigenfunction
with $|E_j-E|=d$; then $c_j=0$ and, by \eqref{expansion},
\[
\bar\Lambda_\mu(u+tu_j,E)=\|(L-E)u\|^2+t^2d^2+\mu(s+t^2-1)^2
=\bar\Lambda_\mu(u,E)+t^2(d^2-\lambda)+\mu t^4,
\]
where we used $2\mu(s-1)=-\lambda$. The right-hand side is smaller than
$\bar\Lambda_\mu(u,E)$ for small $t\neq0$, so $u$ is not a local minimizer. In the same
way, $\bar\Lambda_\mu(tu_j,E)=\mu+t^2(d^2-2\mu)+\mu t^4$, so $u=0$ is not a local minimizer
when $d^2<2\mu$; when $d^2\ge2\mu$ it is the global minimizer by (b). Finally, the critical
points with $\lambda=d^2$ are the global minimizers of (b).
\end{proof}

\begin{remark}[No spurious minima at the continuous level]\label{rem.nospurious}
Proposition \ref{prop.penalized}(c) shows that, before discretization, the loss landscape
is benign: for every $E$ all local minimizers of $\bar\Lambda_\mu(\cdot,E)$ are global, and
every other critical point is a saddle with an explicit descent direction (an
eigenfunction of the eigenvalue closest to $E$). Consequently, the spurious minima that
may be observed in practice are not a feature of the variational problem; they can only
be produced by the three discretization steps discussed in Proposition
\ref{prop.stability} below: the finite network (the parametrization
$\theta\mapsto u_\theta$ is nonlinear, so the landscape in parameter space may have
additional critical points), the finite number of optimization steps, and Monte Carlo
sampling.
\end{remark}

\begin{corollary}[Admissible range of $\mu$]\label{cor.mu}
Let $[E_*,E^*]$ be the scanned window and
\[
D:=\max_{E\in[E_*,E^*]}d(E)
\]
(so $D$ is at most half of the largest gap between consecutive distinct eigenvalues
meeting the window, unless the maximum is attained at an endpoint of the window). Let
$u_E$ be a minimizer of $\bar\Lambda_\mu(\cdot,E)$ and let
$R_\mu(E):=\|(L-E)u_E\|^2$ be the residual term at the minimizer, which is the quantity
recorded during the scan. Then:
\begin{enumerate}
\item[(i)] For every $\mu>0$, $\bar\Lambda_\mu(E)=0$ if and only if $E\in\sigma$, and the
local minima of $E\mapsto\bar\Lambda_\mu(E)$ in the region $\{d(E)^2<2\mu\}$ are exactly
the eigenvalues. Moreover
$0\le d(E)^2-\bar\Lambda_\mu(E)\le d(E)^4/(4\mu)$ if $d(E)^2\le2\mu$.
\item[(ii)] If $\mu\ge D^2$, then for every $E$ in the window $\|u_E\|^2\ge\frac12$ and
\[
R_\mu(E)=d(E)^2\Big(1-\frac{d(E)^2}{2\mu}\Big)\in\big[\tfrac12 d(E)^2,\,d(E)^2\big]
\]
is a strictly increasing function of $d(E)$; hence in the window the zeros and the local
minima of $R_\mu$ are exactly the eigenvalues.
\item[(iii)] If $\frac12D^2\le\mu<D^2$, then $R_\mu$ has spurious local minima at the points of the window where $d$ attains a local maximum larger than
$\sqrt\mu$.
\item[(iv)] If $\mu<\frac12D^2$, then for all $E$ in the window with $d(E)^2\ge2\mu$ the
minimizer is $u_E=0$ and $R_\mu(E)=0$, although $E$ is not an eigenvalue: the residual
curve has spurious zeros and $\bar\Lambda_\mu$ has a plateau at height $\mu$.
\end{enumerate}
\end{corollary}

\begin{proof}
The function $\varphi_\mu(t)=t-t^2/(4\mu)$ is strictly increasing on $[0,2\mu]$, with
$\varphi_\mu(0)=0$, and $\bar\Lambda_\mu(E)=\varphi_\mu(d(E)^2)$ there; since the local
minima of $d$ are the points of $\sigma$, (i) follows from Proposition
\ref{prop.penalized}(a). By Proposition \ref{prop.penalized}(b),
$R_\mu(E)=\psi_\mu(d(E)^2)$ with $\psi_\mu(t)=t(1-t/(2\mu))_+$, which is strictly
increasing on $[0,\mu]$, strictly decreasing on $[\mu,2\mu]$ and vanishes on
$[2\mu,\infty)$; moreover $\|u_E\|^2=1-d(E)^2/(2\mu)$. This gives (ii)--(iv).
\end{proof}

\begin{remark}[A guideline for the choice of $\mu$]\label{rem.mu}
Corollary \ref{cor.mu} and the following three observations lead to a quantitative rule.

\emph{(1) Lower bound from the spectrum.} The condition $\mu\ge D^2$ is necessary and
sufficient for the residual curve of the ideal problem to be free of artifacts, and
$\mu<D^2/2$ produces the collapse $u\approx0$ described in the text. Here $D$ is half of
the largest spectral gap in the window (or the distance from an endpoint of the window to
the nearest eigenvalue, whichever is larger); it can be estimated a priori from Weyl's
law, or a posteriori from a coarse preliminary scan. The value of $\mu$ does \emph{not}
move the minima: by Proposition \ref{prop.penalized}(b) the quotient
$\Lambda(u_E,E)=d(E)^2$ is independent of $\mu$. We therefore recommend recording the
normalized residual $\Lambda(u,E)$ rather than the bare residual term, and discarding the
nodes for which $\|u\|^2<\frac12$.

\emph{(2) Lower bound from the initialization.} Let $w$ be a fixed profile with $\|w\|=1$
and consider the loss along the ray $u=aw$, $a\ge0$:
\[
\bar\Lambda_\mu(aw,E)=a^2q+\mu(a^2-1)^2,\qquad q:=\Lambda(w,E).
\]
If $q\ge2\mu$ this function is increasing in $a$, so a descent method decreases the
amplitude towards the trivial state; if $q<2\mu$ its minimum is attained at
$a^2=1-q/(2\mu)>0$. For the profile $w_0$ of a freshly initialized network (a smooth
function with a few oscillations), Proposition \ref{prop.rayleigh}(a) below gives
\[
q_0(E):=\Lambda(w_0,E)=\big(E-\rho(w_0)\big)^2+\eta(w_0)^2\ \le\ E^2+\|Lw_0\|^2
\qquad(E\ge0,\ L\ge0),
\]
which is the precise form of the heuristic
$\int(\Delta u+Eu)^2\sim c_1E^2+c_2$ used above. Hence, when the network is trained
\emph{from scratch} at the level $E$, the amplitude does not initially collapse provided
$\mu>\frac12q_0(E)$, i.e.\ $\mu\gtrsim\frac12E^2$ for large $E$. For instance, for
$w_0\propto B(x)=1-|x|^2$ in the unit disk one finds $\rho(w_0)=6$ and
$\eta(w_0)^2=12$, so $q_0(E)=(E-6)^2+12$. With a warm start the initial profile is the
minimizer at the previous node, for which $q\approx d(E)^2$, and one is back to the much
weaker condition of item (1). By Proposition \ref{prop.penalized}(c) the trivial state is a
saddle whenever $d(E)^2<2\mu$, so a collapse is not irreversible, but escaping from it is
slow; this explains why the constraint (2) matters for cold starts and periodic
re-initializations but not along a warm-started scan.

\emph{(3) Upper bound from sampling.} If $\hat s_N=\frac1N\sum_{i=1}^Nu(x_i)^2$ is the
Monte Carlo estimate of $s=\|u\|^2$ (normalized measure, i.i.d.\ uniform points), then
\begin{equation}\label{mc.floor}
\mathbb E\big[\mu(\hat s_N-1)^2\big]=\mu(s-1)^2+\frac{\mu}{N}\operatorname{Var}\big(u(x)^2\big),
\end{equation}
so that the sampled penalty adds a noise floor $\mu\operatorname{Var}(u^2)/N$ to the total
loss, and gradient noise of proportional size in the direction of the amplitude
(for a normalized eigenfunction $\operatorname{Var}(u^2)=\int u^4\,dx/|\Omega|-1$ is of
order one; it equals $5/4$ for the eigenfunctions of the square). Thus $\mu$ should not
exceed a moderate multiple of $N\epsilon$, where $\epsilon$ is the detection threshold.

\emph{Rule.} We take
\begin{equation}\label{mu.rule}
\mu=\max\big\{2D^2,\ q_0(E^*)\big\}\quad\text{(or node-wise } \mu(E)=\max\{2D^2,q_0(E)\}\text{)},
\qquad \mu\lesssim N\epsilon,
\end{equation}
where $q_0(E^*)\le (E^*)^2+\|Lw_0\|^2$ is computable at initialization at no cost; when
only warm starts are used, any $\mu\ge 2D^2$ is admissible. The values used in our
experiments lie between the two lower bounds $2D^2$ and $q_0(E^*)$, see Table
\texttt{[MU-TABLE-TBD]}; the empirical sensitivity study reported there
(\texttt{[MU-TABLE-TBD]}: computed eigenvalue, residual and $\|u\|^2$ for
$\mu\in\{\dots\}$) confirms the prediction that the location of the minima is insensitive
to $\mu$ in the range \eqref{mu.rule}, while the collapse $u\approx0$ appears for
$\mu<D^2/2$.
\end{remark}


%% ============================================================================
%% FRAGMENT 2 -- answers R1.6 (optimization / approximation / sampling errors,
%% spurious minima).  Suggested place: new subsection "Stability of the loss
%% curve" after the description of the algorithm.
%% NOTE: the paper calls the number of grid nodes "j"; here it is J.
%% NOTE: \epsilon = detection threshold (paper), \varepsilon = early stopping.
%% ============================================================================

\subsection*{Stability of the detection under training and sampling errors}

The identity $\Lambda(E)=d(E)^2$ concerns exact minimizers. The algorithm, instead,
produces at each node of the grid
\[
E^i=E_*+(i-1)h,\qquad i=1,\dots,J,\qquad E^J=E^*,
\]
a trained network $u^i$ and a number $\tilde\Lambda_i$, the computed value of the
(normalized) loss. The error splits as
\begin{equation}\label{error.split}
\tilde\Lambda_i-\Lambda(E^i)=
\underbrace{\Lambda(u^i,E^i)-\inf_{u\in\mathcal N}\Lambda(u,E^i)}_{\text{optimization}\ \ge0}
+\underbrace{\inf_{u\in\mathcal N}\Lambda(u,E^i)-\Lambda(E^i)}_{\text{approximation}\ \ge0}
+\underbrace{\tilde\Lambda_i-\Lambda(u^i,E^i)}_{\text{sampling}},
\end{equation}
where $\mathcal N\subset D(L)$ is the set of functions $u=B\,\mathrm{NN}_\theta$
representable by the chosen architecture. The first two terms are nonnegative, the third
one has no sign. We do not assume anything on the origin of these errors, only a bound on
their size:
\begin{equation}\label{hyp.delta}
\big|\tilde\Lambda_i-d(E^i)^2\big|\le\delta\qquad\text{for } i=1,\dots,J.
\end{equation}
We say that an interior node $E^m$ ($1<m<J$) is a \emph{discrete local minimum} if
$\tilde\Lambda_m\le\min\{\tilde\Lambda_{m-1},\tilde\Lambda_{m+1}\}$.

\begin{proposition}[Stability of the loss-curve characterization]\label{prop.stability}
Assume \eqref{hyp.delta} and let $\epsilon>0$ be the detection threshold. Set
\[
\rho:=\sqrt{\epsilon+\delta},\qquad \tau:=\sqrt{\tfrac{h^2}{4}+2\delta}\ \le\ \tfrac h2+\sqrt{2\delta}.
\]
\begin{enumerate}
\item[(a)] (No spurious detections away from the spectrum.) If
$\tilde\Lambda_i\le\epsilon$, then $d(E^i)\le\rho$: the interval $[E^i-\rho,E^i+\rho]$
contains an eigenvalue.
\item[(b)] (No missed eigenvalues.) For every eigenvalue $E_k\in[E_*,E^*]$ there is a node
with $|E^i-E_k|\le h/2$ and $\tilde\Lambda_i\le h^2/4+\delta$. Hence, if
$\epsilon\ge h^2/4+\delta$, every eigenvalue in the window is detected by thresholding.
\item[(c)] (Local minima are close to the spectrum.) Every discrete local minimum $E^m$
satisfies
\[
d(E^m)\le\max\Big\{h,\ \frac h2+\frac\delta h\Big\},
\]
irrespective of its value $\tilde\Lambda_m$.
\item[(d)] (Every isolated eigenvalue produces a local minimum.) Let $E_k$ be an
eigenvalue with $[E_k-\tau-h,E_k+\tau+h]\subset[E_*,E^*]$ and
\begin{equation}\label{gap.cond}
\gamma_k:=\dist\big(E_k,\sigma\setminus\{E_k\}\big)>2\tau+h .
\end{equation}
Then there is a discrete local minimum $E^m$ with $|E^m-E_k|\le\tau$ and
$\tilde\Lambda_m\le h^2/4+\delta$; $E_k$ is the only eigenvalue in
$[E^m-\tau,E^m+\tau]$.
\item[(e)] (Exact count.) Assume $\epsilon\ge h^2/4+\delta$ and that any two distinct
eigenvalues at distance at most $\rho$ from the window are more than $4\rho$ apart.
Declare two detected nodes (i.e.\ nodes with $\tilde\Lambda_i\le\epsilon$) equivalent if
their distance is at most $2\rho$. This is an equivalence relation, each class consists of
the detected nodes lying within $\rho$ of one and the same eigenvalue, every eigenvalue
in $[E_*,E^*]$ has a nonempty class, and the node $E^{m}$ minimizing $\tilde\Lambda$ in
the class of $E_k\in[E_*,E^*]$ satisfies $|E^m-E_k|\le\tau$.
\end{enumerate}
If the errors are one-sided, $0\le\tilde\Lambda_i-d(E^i)^2\le\delta$ (no sampling error),
the statements hold with $\rho=\sqrt\epsilon$, $\tau=\sqrt{h^2/4+\delta}$ and with
$\delta/(2h)$ in place of $\delta/h$ in (c).
\end{proposition}

\begin{proof}
Write $d_i=d(E^i)$.

(a) By \eqref{hyp.delta}, $d_i^2\le\tilde\Lambda_i+\delta\le\epsilon+\delta$.

(b) Since $E^1\le E_k\le E^J$ and the grid has step $h$, some node satisfies
$|E^i-E_k|\le h/2$; then $d_i\le h/2$ and
$\tilde\Lambda_i\le d_i^2+\delta\le h^2/4+\delta$.

(c) Let $E_k$ be an eigenvalue with $|E^m-E_k|=d_m$ and assume $d_m\ge h$ (otherwise
there is nothing to prove). Let $E^{m'}$, $m'\in\{m-1,m+1\}$, be the neighbour of $E^m$
on the side of $E_k$. Then $d_{m'}\le|E^{m'}-E_k|=d_m-h$, and by \eqref{hyp.delta} and
the definition of discrete local minimum,
\[
d_m^2-\delta\le\tilde\Lambda_m\le\tilde\Lambda_{m'}\le(d_m-h)^2+\delta ,
\]
that is, $2hd_m\le h^2+2\delta$.

(d) Let $I=[E_k-\tau-h,E_k+\tau+h]$ and let $E^m$ minimize $\tilde\Lambda$ among the nodes
in $I$. By (b), $\tilde\Lambda_m\le h^2/4+\delta$, hence $d_m^2\le h^2/4+2\delta=\tau^2$.
If the eigenvalue closest to $E^m$ were some $E_j\neq E_k$, we would have
$|E^m-E_k|\ge\gamma_k-\tau>\tau+h$ by \eqref{gap.cond}, contradicting $E^m\in I$.
Therefore $|E^m-E_k|=d_m\le\tau$; in particular the two neighbours of $E^m$ belong to
$I$, so $\tilde\Lambda_m\le\tilde\Lambda_{m\pm1}$, and $E^m$ is an interior node because
$I\subset[E_*,E^*]$. Any other eigenvalue is at distance at least
$\gamma_k-\tau>\tau$ from $E^m$.

(e) By (a) every detected node lies within $\rho$ of some eigenvalue, which is then at
distance at most $\rho$ from the window and is unique by the separation assumption. Two
detected nodes attached to the same eigenvalue are at distance at most $2\rho$; two
detected nodes attached to different eigenvalues are at distance larger than
$4\rho-2\rho=2\rho$. This proves that the relation is an equivalence relation whose
classes are indexed by eigenvalues. By (b) the class of an eigenvalue $E_k$ in the window
contains a node with $\tilde\Lambda_i\le h^2/4+\delta$; thus the minimizing node of the
class satisfies $d_m^2\le h^2/4+2\delta$ and, since it is attached to $E_k$,
$|E^m-E_k|=d_m\le\tau$.

The one-sided case follows by repeating the arguments with $d_i^2\le\tilde\Lambda_i$ in
place of $d_i^2-\delta\le\tilde\Lambda_i$.
\end{proof}

\begin{remark}[Interpretation]\label{rem.stability}
(1) \emph{Spurious minima.} Imperfect training cannot create detections far from the
spectrum: by (a) and (c) every node below the threshold, and every discrete local minimum
whatsoever, lies within an explicit distance of a true eigenvalue. What imperfect
training \emph{can} do is (i) produce several discrete local minima around the same
eigenvalue, all of them within $\max\{h,h/2+\delta/h\}$ of it, which are merged by the
grouping in (e); and (ii) raise the whole curve above the threshold, so that an
eigenvalue is missed if $\epsilon<h^2/4+\delta$. The bound in (c) also shows that refining
the grid below $h\simeq\sqrt{2\delta}$ does not improve the localization: the quantity
$h/2+\delta/h$ is minimal, and equal to $\sqrt{2\delta}$, for $h=\sqrt{2\delta}$.

(2) \emph{Resolution.} The eigenvalue is located within $\tau\le h/2+\sqrt{2\delta}$ by the
grid alone; the sub-grid procedures of Propositions \ref{prop.rayleigh} and
\ref{prop.parabola} improve on this.

(3) \emph{Multiplicity.} The loss curve only sees the set $\sigma$: a multiple eigenvalue
produces a single minimum, and the trained network converges to \emph{some} element of
the eigenspace. Multiplicities have to be recovered separately, for instance by training
at the detected value from several independent initializations and computing the
numerical rank of the Gram matrix of the resulting functions.

(4) \emph{Clusters of eigenvalues.} Two eigenvalues $E_k<E_{k+1}$ are separated by a bump
of the loss curve of height $(E_{k+1}-E_k)^2/4$. If $E_{k+1}-E_k\le2\sqrt{2\delta}$
(respectively $2\sqrt\delta$ for one-sided errors) this bump is below the error level:
the constant data $\tilde\Lambda\equiv\delta$ (respectively
$\tilde\Lambda\equiv(E_{k+1}-E_k)^2/4$) on the nodes of $[E_k,E_{k+1}]$ are compatible
with \eqref{hyp.delta} for the true spectrum and also for any spectrum containing
additional eigenvalues in $(E_k,E_{k+1})$. Hence the number of eigenvalues in a cluster
of width below $2\sqrt{2\delta}$ cannot be determined from the loss curve; the method then
returns the cluster as a single approximate eigenvalue, with the error bound of (a).

(5) \emph{What is not proved.} Proposition \ref{prop.stability} is conditional on
\eqref{hyp.delta}. We do not prove a bound for $\delta$ in terms of the width and depth of
the network, the number of epochs and the number of sampling points. The approximation
term in \eqref{error.split} tends to zero as the network grows (Remark
\ref{rem.approx}), but quantitative rates for the minimizers at hand, and above all a
convergence guarantee for the optimizer, are not available; in particular the spectral
bias of gradient-based training \cite{Rahaman-etal} suggests that $\delta$ grows with the
oscillation of the eigenfunction, i.e.\ with $E$, in agreement with the residuals in our
tables. What can always be done is to \emph{measure} the error a posteriori, see
Proposition \ref{prop.rayleigh}.
\end{remark}

\begin{remark}[Consistency of the network class]\label{rem.approx}
Assume that the eigenfunction $u_k$ can be written as $u_k=Bv_k$ with
$v_k\in C^2(\overline\Omega)$ and $B\in W^{2,\infty}(\Omega)$ (this holds, for instance,
when $\partial\Omega$ and $V$ are smooth and $B$ is smooth with $\nabla B\neq0$ on
$\partial\Omega$, and for the explicit eigenfunctions and polynomial factors $B$ of the
square and of the right isosceles triangle). Since networks with one hidden layer and
smooth non-polynomial activation are dense in $C^2(\overline\Omega)$
\cite{Hornik1991,Pinkus1999}, there are networks $\mathrm{NN}_n\to v_k$ in
$C^2(\overline\Omega)$; then $B\,\mathrm{NN}_n\to u_k$ in $H^2(\Omega)$, and by continuity
of $\Lambda(\cdot,E)$ in $H^2(\Omega)\setminus\{0\}$,
\[
\inf_{\text{all widths}}\ \Lambda(B\,\mathrm{NN}_\theta,E)=\Lambda(E)\qquad\text{for all }E .
\]
For $\tanh$ networks explicit rates in Sobolev norms in terms of the width are given in
\cite{DeRyck-Lanthaler-Mishra}; convergence of the minimizers of sampled PINN losses for
linear elliptic source problems is studied in \cite{Shin-Darbon-Karniadakis}. Neither
result covers the optimization error.
\end{remark}


%% ============================================================================
%% FRAGMENT 3 -- answers R1.6 (a posteriori certificate) and explains
%% "small eigenvalue error with large residual" rigorously.  Suggested place:
%% replace the paragraph "Eigenvalue accuracy versus PDE residual".
%% ============================================================================

\subsection*{A posteriori certificates and the role of the Rayleigh quotient}

For $u\in D(L)\setminus\{0\}$ let
\[
\rho(u):=\frac{\langle Lu,u\rangle}{\|u\|^2},\qquad
\eta(u)^2:=\frac{\|Lu-\rho(u)u\|^2}{\|u\|^2}=\frac{\|Lu\|^2}{\|u\|^2}-\rho(u)^2 .
\]

\begin{proposition}\label{prop.rayleigh}
Let $u\in D(L)\setminus\{0\}$ be arbitrary (for instance, a trained network).
\begin{enumerate}
\item[(a)] For every $E\in\R$,
\[
\Lambda(u,E)=\big(E-\rho(u)\big)^2+\eta(u)^2 .
\]
In particular, for fixed $u$, the function $E\mapsto\Lambda(u,E)$ is a parabola with
leading coefficient $1$, whose minimum is attained at the Rayleigh quotient $\rho(u)$
and equals $\eta(u)^2$.
\item[(b)] (Weinstein.) For every $E\in\R$, $d(E)\le\Lambda(u,E)^{1/2}$. In particular
$d(\rho(u))\le\eta(u)$.
\item[(c)] (Kato--Temple \cite{Kato1949,Parlett1998}.) If $\alpha<\rho(u)<\beta$ and $E_k$ is the only point of
$\sigma$ in $(\alpha,\beta)$, then
\[
\rho(u)-\frac{\eta(u)^2}{\beta-\rho(u)}\ \le\ E_k\ \le\ \rho(u)+\frac{\eta(u)^2}{\rho(u)-\alpha}.
\]
In particular, if $E_k$ is the only point of $\sigma$ in
$(\rho(u)-\gamma,\rho(u)+\gamma)$, then
\[
|\rho(u)-E_k|\le\frac{\eta(u)^2}{\gamma}\le\frac{\Lambda(u,E)}{\gamma}\qquad\text{for every }E\in\R.
\]
\end{enumerate}
\end{proposition}

\begin{proof}
We may assume $\|u\|=1$; write $u=\sum_jc_ju_j$ and $\rho=\rho(u)$, $\eta=\eta(u)$.

(a) $\|(L-E)u\|^2=\|(L-\rho)u\|^2+2(\rho-E)\langle(L-\rho)u,u\rangle+(\rho-E)^2$, and
$\langle(L-\rho)u,u\rangle=0$.

(b) This is \eqref{expansion}: $\|(L-E)u\|^2=\sum_j(E_j-E)^2c_j^2\ge d(E)^2$.

(c) Since $\sigma\cap(E_k,\beta)=\emptyset$, we have $(E_j-E_k)(E_j-\beta)\ge0$ for all $j$.
Writing $E_j-E_k=(E_j-\rho)+(\rho-E_k)$, $E_j-\beta=(E_j-\rho)+(\rho-\beta)$ and using
$\sum_jc_j^2=1$, $\sum_jc_j^2(E_j-\rho)=0$ and $\sum_jc_j^2(E_j-\rho)^2=\eta^2$, we get
\[
0\le\sum_jc_j^2(E_j-E_k)(E_j-\beta)=\eta^2+(\rho-E_k)(\rho-\beta),
\]
that is, $(\rho-E_k)(\beta-\rho)\le\eta^2$, which is the first inequality. The second one
follows in the same way from $(E_j-\alpha)(E_j-E_k)\ge0$. The last assertion follows by
taking $\alpha=\rho-\gamma$, $\beta=\rho+\gamma$, and from (a).
\end{proof}

\begin{remark}[Consequences for the method]\label{rem.certificate}
(1) \emph{Certificate independent of the training.} Items (b) and (c) hold for
\emph{any} $u\in D(L)$. Since $u=B\,\mathrm{NN}$ satisfies the boundary condition exactly
and is smooth, the trained network is an admissible $u$, whatever the quality of the
training. Hence each detected pair $(u^i,E^i)$ comes with the guaranteed enclosure
\[
\sigma\cap\big[E^i-r_i,\,E^i+r_i\big]\neq\emptyset,\qquad r_i:=\Lambda(u^i,E^i)^{1/2},
\]
which turns a possibly spurious minimum into a verified one: a detection is accepted
when $r_i$ is small compared with the distance to the neighbouring detections. The
quantity $\Lambda(u^i,E^i)$ is an integral that can be evaluated after the training with
a fresh and arbitrarily large sample, see (4). (The column ``Residual'' of our tables
reports $\Lambda(u^i,E^i)=r_i^2$.)

(2) \emph{Local refinement is a Rayleigh quotient.} Once a node $E^{i_0}$ has been
detected, minimizing $E\mapsto\Lambda(u^{i_0},E)$ with the network frozen needs no
further training and, by (a), returns exactly $\rho(u^{i_0})$, which can be computed in
closed form:
\[
\hat E=\rho(u^{i_0})=\frac{\int_\Omega u^{i_0}\,Lu^{i_0}\,dx}{\int_\Omega (u^{i_0})^2\,dx},
\qquad \eta(u^{i_0})^2=\Lambda(u^{i_0},\hat E).
\]

(3) \emph{Why the eigenvalue is more accurate than the residual.} By (c), if $\gamma$ is
(a lower bound for) the distance from $\hat E$ to the rest of the spectrum, then
\[
|\hat E-E_k|\le\frac{\eta(u^{i_0})^2}{\gamma},
\]
whereas the residual norm is $\eta(u^{i_0})$: the eigenvalue error is \emph{quadratic}
in the residual. For example, a residual $\eta^2=10^{-1}$ and a gap $\gamma=10$ guarantee
$|\hat E-E_k|\le10^{-2}$ although $\eta\approx0.32$. This replaces the heuristic argument
of a ``constant floor'' and explains quantitatively the pattern observed in all the
tables: accurate eigenvalues together with comparatively large residuals. The gap
$\gamma$ can be bounded from below using the enclosures of the neighbouring detections
given by (b), so the bound is computable.

(4) \emph{Monte Carlo evaluation, variance and bias.} Let $x_1,\dots,x_N$ be i.i.d.\
uniform in $\Omega$ and independent of $u$ (fresh points), and let
$g=(Lu-Eu)^2$, $\hat R_N=\frac1N\sum_ig(x_i)$, $\hat s_N=\frac1N\sum_iu(x_i)^2$. Each
average is an unbiased estimator of the corresponding (normalized) integral, with
variance $\operatorname{Var}(g)/N$ and $\operatorname{Var}(u^2)/N$. The loss is, however, a
nonlinear function of these averages, and this introduces a bias of order $1/N$:
for the penalty term the bias is exactly the one in \eqref{mc.floor}, and for the
quotient $\hat\Lambda_N=\hat R_N/\hat s_N$ the delta method gives, with $R=\mathbb E g$,
$s=\mathbb E u^2$, $\Lambda=R/s$,
\[
\mathbb E\hat\Lambda_N-\Lambda=\frac1N\Big(\Lambda\,\frac{\operatorname{Var}(u^2)}{s^2}
-\frac{\operatorname{Cov}(g,u^2)}{s^2}\Big)+O(N^{-2}),
\qquad
\operatorname{Var}\hat\Lambda_N=\frac{\operatorname{Var}(g-\Lambda u^2)}{Ns^2}+O(N^{-2}).
\]
Thus the bias ($O(1/N)$) is negligible with respect to the statistical error
($O(N^{-1/2})$), and the latter can be estimated from the same sample: with
$\hat\sigma_N^2$ the empirical variance of $g(x_i)-\hat\Lambda_Nu(x_i)^2$,
\[
d(E)\le\Big(\hat\Lambda_N+z\,\frac{\hat\sigma_N}{\hat s_N\sqrt N}\Big)^{1/2}
\]
holds with asymptotic confidence $\Phi(z)$ (e.g.\ $97.7\%$ for $z=2$). This is a
statistical certificate, not a deterministic one. The same remark applies to
$\rho(u)$, which is also a ratio of two averages. Two further points deserve mention.
First, during the training the sample is redrawn at every epoch, and the value recorded
at the last epoch is computed on a sample that has not yet been used to update the
parameters; it is therefore an unbiased evaluation of the residual term for the current
network, and no optimistic bias due to a fixed training sample is present. Second, the
sub-grid refinement must not be performed by minimizing a \emph{re-sampled} objective
$E\mapsto\hat\Lambda_N(u,E)$: by (a) this objective is flat to within its noise level
$\nu\sim\hat\sigma_N/\sqrt N$ on an interval of half-width $\sqrt\nu$ around $\rho(u)$,
so the returned minimizer has an error of order $\sqrt\nu$, much larger than the error
$O(N^{-1/2})$ of the closed formula in (2) evaluated on a single sample.
\end{remark}


%% ============================================================================
%% FRAGMENT 4 -- sub-grid accuracy by parabolic interpolation (R1.1, R1.6).
%% Suggested place: cost subsection, or right after Fragment 3.
%% ============================================================================

\begin{proposition}[Three-point parabolic interpolation]\label{prop.parabola}
Let $E^{m-1},E^m,E^{m+1}$ be three consecutive nodes and $y_-,y_0,y_+$ the corresponding
computed values of the loss. Define
\begin{equation}\label{parabola.vertex}
\hat E:=E^m+\frac h2\,\frac{y_--y_+}{y_--2y_0+y_+},\qquad
\hat a:=\frac{y_--2y_0+y_+}{2h^2}.
\end{equation}
Assume that the three nodes lie in $\{E\colon|E-E_k|=d(E)\}$, i.e.\ $E_k$ is the
eigenvalue closest to each of them, and that for some constant $c\in\R$
\[
\big|y_\bullet-(E^{\bullet}-E_k)^2-c\big|\le\omega\qquad\text{at the three nodes, with } 2\omega<h^2 .
\]
Then, with $e:=E^m-E_k$,
\[
|\hat E-E_k|\le\frac{(4|e|+h)\,\omega}{2h^2-4\omega},\qquad|\hat a-1|\le\frac{2\omega}{h^2}.
\]
In particular, if $E^m$ is the node closest to $E_k$ ($|e|\le h/2$) and
$\omega\le h^2/4$, then $|\hat E-E_k|\le3\omega/h$.
\end{proposition}

\begin{proof}
Let $\varepsilon_\bullet=y_\bullet-(E^\bullet-E_k)^2-c$, so $|\varepsilon_\bullet|\le\omega$.
Since $(e\mp h)^2$ are the exact values at $E^{m\mp1}$,
\[
y_--y_+=-4he+n,\quad y_--2y_0+y_+=2h^2+m,\qquad
n:=\varepsilon_--\varepsilon_+,\quad m:=\varepsilon_--2\varepsilon_0+\varepsilon_+,
\]
with $|n|\le2\omega$, $|m|\le4\omega$ (the constant $c$ cancels). Therefore
\[
\hat E-E_k=e+\frac h2\,\frac{-4he+n}{2h^2+m}=\frac{em+hn/2}{2h^2+m},
\]
and the estimates follow from $2h^2+m\ge2h^2-4\omega>0$.
\end{proof}

\begin{remark}\label{rem.parabola}
For exact data $y=(E-E_k)^2+c$ the formula \eqref{parabola.vertex} returns $E_k$ exactly,
for any $h$: a constant error floor $c$ (the same excess loss at the three nodes) does
not affect the result, only its oscillation $\omega$ does. With the grid alone the error
is $h/2$; with three nodes it is $3\omega/h$, which is smaller as soon as
$\omega<h^2/6$. This allows the use of coarse grids. The number $\hat a$ is a useful
diagnostic: in the linear case it must be close to $1$ (for the unnormalized loss
$\bar\Lambda_\mu$ the leading coefficient is $1-3d^2/(2\mu)$ up to $O(h^2/\mu)$, by
Proposition \ref{prop.penalized}(a)), and a value far from $1$ reveals that the three nodes do not belong to the same parabola, i.e.\ that the
training has not converged at some of them or that two eigenvalues are involved.
\end{remark}


%% ============================================================================
%% FRAGMENT 5 -- answers R1.1 (cost of the scan, comparison with Rayleigh PINNs)
%% Suggested place: new subsection "Computational cost".
%% ============================================================================

\subsection*{Computational cost of the scan}

\emph{Cost model.} Let $W$ be the number of parameters of the network, $N$ the number of
sampling points and $d$ the space dimension. One evaluation of $u_\theta$ at the $N$
points costs $c_0NW$ operations. The loss \eqref{lossmu} requires $\Delta u_\theta$,
which by reverse-mode automatic differentiation costs one backward pass for
$\nabla u_\theta$ and $d$ further backward passes for the diagonal of the Hessian;
differentiating the result with respect to $\theta$ multiplies the cost by a constant.
Hence one epoch costs
\[
c_{\rm ep}=\kappa\,(d+2)\,NW
\]
with $\kappa$ independent of $d$, $N$ and $W$: the cost per epoch grows linearly, not
exponentially, with the dimension. If $n_i$ denotes the number of epochs spent at the node
$E^i$, the cost of the scan of $[E_*,E^*]$ with $J=1+(E^*-E_*)/h$ nodes is
\begin{equation}\label{cost.scan}
C_{\rm scan}=c_{\rm ep}\sum_{i=1}^Jn_i\ \le\ \kappa\,(d+2)\,N\,W\,J\,n_{\max},
\end{equation}
where $n_{\max}$ is the maximal number of epochs. Three features reduce this bound.
\begin{enumerate}
\item \emph{Warm start.} The minimizer at $E^{i-1}$ is an eigenfunction of the eigenvalue
closest to $E^{i-1}$, which is also the minimizer at $E^i$ unless the midpoint between
two eigenvalues lies in $(E^{i-1},E^i]$. Thus, with warm start and early stopping, only
the nodes following such a midpoint (at most one per eigenvalue) require a full
training, and $\sum_in_i\approx K_In_{\max}+(J-K_I)\,\bar n$, where $K_I$ is the number
of distinct eigenvalues in the window and $\bar n\ll n_{\max}$ is the average number of
epochs at the remaining nodes; in our experiments $\bar n/n_{\max}\approx$
\texttt{[RUNTIME-TBD]}.
\item \emph{Coarse grids with sub-grid refinement.} By Proposition \ref{prop.stability}
the grid only has to \emph{separate} the eigenvalues: it suffices that
$h<\gamma_{\min}-2\tau$, i.e.\ essentially $h\lesssim\gamma_{\min}/2$, where
$\gamma_{\min}$ is the smallest gap between distinct eigenvalues in the window. The
accuracy is then obtained without further training from the Rayleigh quotient of the
detected network (Remark \ref{rem.certificate}(2)) or from the three-point formula
\eqref{parabola.vertex}, and it is not limited by $h$. Hence
$J\approx2(E^*-E_*)/\gamma_{\min}$, that is, a small number of nodes per eigenvalue when
the eigenvalues in the window are well separated. If only an accuracy $h$ were sought by
refining the grid, the cost would be proportional to $1/h$; an adaptive strategy (coarse
scan, then local grids around the detected nodes) reduces this to
$O(K_I\log(1/h))$ trainings by bisection, and to $O(K_I)$ with the formulas above.
\item \emph{Parallelism.} The trainings at different values of $E$ are independent. With
$P$ processors the window is split into $P$ subwindows scanned simultaneously (each one
with its own warm start), and the wall-clock time is $C_{\rm scan}/P$ up to the cost of
$P-1$ additional cold starts. Measured times: \texttt{[RUNTIME-TBD]}.
\end{enumerate}

\emph{Comparison with Rayleigh-quotient PINNs.} In the methods based on the Rayleigh
quotient described in the Introduction (see
\cite{Rowan-Evans-Maute-Dootsan,Holliday-Lindner-Ditto}, and \cite{E-Yu} for the Deep Ritz
formulation) the $k$-th eigenpair is obtained by minimizing $L(u)$ plus penalties enforcing orthogonality
to the $k-1$ previously computed networks. One epoch needs only first derivatives of
$u_\theta$ (cost $\kappa'NW$, independent of $d$), plus the evaluation of the previous
$k-1$ networks at the sampling points (cost $c_0(k-1)NW$). Therefore the computation of
the first $K$ eigenpairs costs
\begin{equation}\label{cost.rq}
C_{\rm RQ}(K)=\sum_{k=1}^Kn_k\big(\kappa'+c_0(k-1)\big)NW
\ \approx\ \bar n_{\rm RQ}\,NW\Big(\kappa'K+\frac{c_0}{2}K^2\Big),
\end{equation}
and it is intrinsically sequential: the $k$-th training cannot start before the previous
ones have finished, and its result depends on their accuracy. Indeed, by the
Courant--Fischer principle,
\[
\min\big\{L(u)\colon u\perp\tilde u_1,\dots,\tilde u_{k-1}\big\}\le E_k
\]
for \emph{any} functions $\tilde u_1,\dots,\tilde u_{k-1}$, with equality when they span
the first $k-1$ eigenfunctions; hence, already for exactly enforced constraints, errors
in the computed eigenfunctions bias all the subsequent eigenvalues downwards and
accumulate with $k$. In the scan, on the contrary,
each node is an independent problem and comes with its own a posteriori bound
(Proposition \ref{prop.rayleigh}).

Comparing \eqref{cost.scan} and \eqref{cost.rq} we obtain the following picture. Let
$K_-$ be the number of eigenvalues below $E_*$ and $K_I$ the number of eigenvalues in the
window, counted with multiplicity. The deflation approach needs $K_-+K_I$ trainings and
a cost growing quadratically in $K_-+K_I$, while the scan needs $J\approx cK_I$ nodes
($c$ a small constant depending on the ratio between the average and the minimal gap),
of which only about $K_I$ are full trainings, at a cost per epoch larger by the factor
$\kappa(d+2)/\kappa'$ due to the second derivatives:
\[
\frac{C_{\rm scan}}{C_{\rm RQ}}\approx
\frac{\kappa(d+2)\,\big(K_In_{\max}+(J-K_I)\bar n\big)}
{\bar n_{\rm RQ}\big(\kappa'(K_-+K_I)+\tfrac{c_0}{2}(K_-+K_I)^2\big)} .
\]
Consequently the scan is \emph{not} competitive for the lowest few eigenvalues
($K_-=0$, $K_I$ small), for which the Rayleigh quotient is cheaper per epoch and needs no
grid; it becomes advantageous for windows located high in the spectrum, since by Weyl's
law $K_-\sim c_d|\Omega|E_*^{d/2}$ grows with $E_*$ and with the dimension while the cost
of the scan depends only on the window, and whenever parallel resources are available.
For the examples of this paper, the measured wall-clock times are
\texttt{[RUNTIME-TBD]} (scan, $J=\ $\texttt{[RUNTIME-TBD]} nodes) versus
\texttt{[RUNTIME-TBD]} (Rayleigh quotient with deflation, same network and sample size).
We stress that the scan does not provide multiplicities (Remark
\ref{rem.stability}(3)), which the deflation approach does, and that it is the only one
of the two that applies to the nonlinear problems of Section \ref{sec.nolin}, where no
orthogonality is available.


%% ============================================================================
%% FRAGMENT 6 -- answers R1.5 (p-Laplacian: error sources, what is and is not
%% proved).  Suggested place: Section "Nonlinear problems", after the definition
%% of Lambda_p, replacing the paragraph "As in the linear case, define ...".
%% ============================================================================

\subsection*{The loss curve for the $p$-Laplacian: what can be proved}

Let $\Omega\subset\R^d$ be a bounded domain, $1<p<\infty$, and write
$\Delta_pu=\nabla\cdot(|\nabla u|^{p-2}\nabla u)$,
$\mathcal R_E(u):=\Delta_pu+E|u|^{p-2}u$. We say that $E$ is an eigenvalue if
\eqref{p-eigen} has a weak solution $u\in W^{1,p}_0(\Omega)\setminus\{0\}$, and denote by
$\sigma_p$ the set of eigenvalues. We recall that $\sigma_p\subset(0,\infty)$ is closed,
that the first eigenvalue
$E_1=\min\{\int_\Omega|\nabla u|^p\,dx\colon \|u\|_p=1\}$ is simple and isolated
\cite{Anane1987,Lindqvist1990}, that there is a sequence of variational eigenvalues given
by the Lusternik--Schnirelmann minimax principle \cite{GarciaAzorero-Peral}, and that the
second eigenvalue $E_2=\min(\sigma_p\setminus\{E_1\})$ coincides with the second
variational eigenvalue \cite{Anane-Tsouli} and admits a mountain-pass characterization
\cite{Cuesta-deFigueiredo-Gossez}, which is the basis of the algorithm of \cite{HJ11}.
It is not known whether the variational eigenvalues exhaust $\sigma_p$. Eigenfunctions
are bounded and of class $C^{1,\alpha}$ \cite{DiBenedetto,Tolksdorf}, but in general not
of class $C^2$.

The natural class for the residual loss is
\[
\mathcal D_p:=\big\{u\in W^{1,p}_0(\Omega)\cap L^{2(p-1)}(\Omega)\colon
\Delta_pu\in L^2(\Omega)\big\},
\]
where $\Delta_pu$ is understood in the sense of distributions. Every eigenfunction belongs
to $\mathcal D_p$, since it is bounded and $\Delta_pu=-E|u|^{p-2}u$.

\begin{proposition}\label{prop.plap}
Let $p\ge\frac{2d}{d+2}$, $p>1$ (so that $W^{1,p}_0(\Omega)\subset L^2(\Omega)$; this is
no restriction if $d=2$).
\begin{enumerate}
\item[(a)] (Homogeneity.) For $t\in\R$ and $u\in\mathcal D_p$,
$\mathcal R_E(tu)=|t|^{p-2}t\,\mathcal R_E(u)$. Hence, if $u$ is an eigenfunction so is
$tu$, with the same eigenvalue, and the quotient
\[
Q_p(u,E):=\frac{\int_\Omega\mathcal R_E(u)^2\,dx}{\big(\int_\Omega|u|^p\,dx\big)^{2(p-1)/p}}
\]
is homogeneous of degree zero. Let $q_*(E):=\inf\{Q_p(u,E)\colon u\in\mathcal D_p\setminus\{0\}\}$.
\item[(b)] (Reduction of the penalized loss.) For $\mu>0$,
\[
\Lambda_p(E)=\inf_{u\in\mathcal D_p}\Lambda_p(u,E)=\Psi_\mu\big(q_*(E)\big),\qquad
\Psi_\mu(q):=\min_{t\ge0}\Big[t^{2(p-1)}q+\mu(t^p-1)^2\Big].
\]
The function $\Psi_\mu$ is nondecreasing, $\Psi_\mu(0)=0$, and
$0<\Psi_\mu(q)\le\min\{q,\mu\}$ for $q>0$.
\item[(c)] (Upper bound.) If $u_k$ is an eigenfunction with eigenvalue $E_k$ and
$\|u_k\|_p=1$, then $Q_p(u_k,E)=c_k(E-E_k)^2$ with
$c_k=\int_\Omega|u_k|^{2(p-1)}dx$. Hence
\[
\Lambda_p(E)\le q_*(E)\le\inf_kc_k(E-E_k)^2 ,
\]
the infimum being taken over all normalized eigenfunctions.
\item[(d)] (The zeros of the loss curve are exactly the eigenvalues.)
$q_*(E)=0$ if and only if $E\in\sigma_p$. Equivalently,
\[
\Lambda_p(E)=0\iff E\in\sigma_p ,
\]
and $\Lambda_p(E)>0$ for every $E\notin\sigma_p$.
\end{enumerate}
\end{proposition}

\begin{proof}
(a) is immediate from $\Delta_p(tu)=|t|^{p-2}t\Delta_pu$.

(b) Every $u\in\mathcal D_p\setminus\{0\}$ can be written as $u=tw$ with $t=\|u\|_p>0$
and $\|w\|_p=1$, and by (a)
$\Lambda_p(tw,E)=t^{2(p-1)}Q_p(w,E)+\mu(t^p-1)^2$, while $\Lambda_p(0,E)=\mu$, which is
the value of the same expression at $t=0$. For each $t\ge0$ the expression is
nondecreasing and affine in $Q_p(w,E)$, so taking first the infimum over $w$ and then
over $t$ gives the formula. The choices $t=1$ and $t=0$ give
$\Psi_\mu(q)\le\min\{q,\mu\}$; if $q>0$, the continuous function
$t\mapsto t^{2(p-1)}q+\mu(t^p-1)^2$ is positive on $[0,\infty)$ and tends to infinity, so
its minimum is positive.

(c) Since $\Delta_pu_k=-E_k|u_k|^{p-2}u_k$, we have
$\mathcal R_E(u_k)=(E-E_k)|u_k|^{p-2}u_k$. The first inequality is contained in (b).

(d) If $E\in\sigma_p$ then $q_*(E)=0$ by (c). Conversely, let $q_*(E)=0$ and take
$w_n\in\mathcal D_p$ with $\|w_n\|_p=1$ and
$f_n:=-\mathcal R_E(w_n)=-\Delta_pw_n-E|w_n|^{p-2}w_n\to0$ in $L^2(\Omega)$. Testing the
equation $-\Delta_pw_n=E|w_n|^{p-2}w_n+f_n$ with $w_n\in W^{1,p}_0(\Omega)\subset
L^2(\Omega)$ and using the Sobolev inequality $\|w\|_2\le C_S\|\nabla w\|_p$,
\[
\|\nabla w_n\|_p^p=E+\int_\Omega f_nw_n\,dx\le E+C_S\|f_n\|_2\|\nabla w_n\|_p ,
\]
so $\{w_n\}$ is bounded in $W^{1,p}_0(\Omega)$. Up to a subsequence, $w_n\rightharpoonup w$
weakly in $W^{1,p}_0(\Omega)$ and strongly in $L^p(\Omega)$; thus $\|w\|_p=1$ and
$|w_n|^{p-2}w_n\to|w|^{p-2}w$ in $L^{p'}(\Omega)$, hence in $W^{-1,p'}(\Omega)$. Moreover
$f_n\to0$ in $L^2(\Omega)\hookrightarrow W^{-1,p'}(\Omega)$ (the adjoint of the Sobolev
embedding). Therefore $-\Delta_pw_n\to E|w|^{p-2}w$ in $W^{-1,p'}(\Omega)$, and since
$(-\Delta_p)^{-1}\colon W^{-1,p'}(\Omega)\to W^{1,p}_0(\Omega)$ is continuous,
$w_n$ converges strongly in $W^{1,p}_0(\Omega)$; the limit is $w$ and it satisfies
$-\Delta_pw=E|w|^{p-2}w$ with $w\neq0$. Thus $E\in\sigma_p$. The equivalence with
$\Lambda_p(E)=0$ follows from (b).
\end{proof}

\begin{remark}[What is lost with respect to the linear case]\label{rem.plap.lost}
Proposition \ref{prop.plap}(d) is the rigorous justification of the method in the
nonlinear case: the loss curve vanishes exactly on the spectrum, no orthogonality being
needed. However, the three quantitative tools of the linear theory are missing.

(1) \emph{No lower bound.} The proof of (d) is by compactness and gives no estimate of
the form $q_*(E)\ge c\,\dist(E,\sigma_p)^2$. Consequently there is no analogue of the
Weinstein bound (Proposition \ref{prop.rayleigh}(b)): a small value of the loss at $E$
does not come with a computable bound on $\dist(E,\sigma_p)$, the stability result
(Proposition \ref{prop.stability}) is not available, and a low minimum of the computed
curve for which no independent value is known (such as the third row of Table
\ref{table.5}) must be regarded as a candidate, not as a certified eigenvalue. Likewise,
only the zeros, and not the positive local minima, of $\Lambda_p$ are characterized, and
the existence of a minimizer for $E\notin\sigma_p$ is not known.

(2) \emph{The curve is not a universal parabola.} The constants $c_k$ in (c) depend on
the eigenfunction and on $p$ (for $p\ge2$, $c_k\le\|u_k\|_\infty^{p-2}$), and the upper
bound need not be attained; hence the curvature of the wells is not known a priori and
the diagnostic $\hat a\approx1$ of Remark \ref{rem.parabola} is lost.

(3) \emph{The refined eigenvalue is only first-order accurate.} For a fixed function $u$
the map $E\mapsto\int\mathcal R_E(u)^2dx$ is still a parabola, minimized at
\[
\rho_p^*(u)=\frac{-\int_\Omega\Delta_pu\,|u|^{p-2}u\,dx}{\int_\Omega|u|^{2(p-1)}dx},
\]
which is the value returned by the local refinement. For $p=2$ this is the Rayleigh
quotient, whose error is quadratic in the error of the eigenfunction (Proposition
\ref{prop.rayleigh}(c)). For $p\neq2$, $\rho_p^*$ is \emph{not} the variational quotient
$\mathcal R_p(u)=\int|\nabla u|^p/\int|u|^p$: eigenfunctions are critical points of
$\mathcal R_p$ but in general not of $\rho_p^*$, so the error of $\rho_p^*(u)$ is
generically of first order in $u-u_k$. This is, in our opinion, the main reason for the
loss of accuracy observed in Tables \ref{table.5} and \ref{table.6} with respect to the
linear examples at comparable residuals. A natural remedy, which we adopt in
\texttt{[PLAP-TBD]}, is to report also $\mathcal R_p(u)$ for the trained network; for
the first eigenvalue this number is moreover a guaranteed upper bound,
$E_1\le\mathcal R_p(u)$ for every $u\in W^{1,p}_0(\Omega)\setminus\{0\}$.

(4) \emph{Role of the normalization.} Because of (a), the equation determines $u$ only up
to a multiplicative constant, and the penalty $\mu(\int|u|^p-1)^2$ only fixes this
constant; by (b), it does not move the zeros of the loss curve. It is essential, however,
that the exponents in the equation and in the constraint be the homogeneous ones: for
the non-homogeneous problem $-\Delta_pu=Eu$ the function $tu$ solves the equation with
$E$ replaced by $|t|^{p-2}E$, so that the computed ``eigenvalue'' would depend on the
normalization chosen and on $|\Omega|$.
\end{remark}

\begin{lemma}[Degeneracy at critical points]\label{lem.critical}
Let $p>2$, $u\in C^2$ in a neighbourhood of $x_0$ and $\nabla u(x_0)=0$. Then
$|\nabla u|^{p-2}\nabla u$ is of class $C^1$ near $x_0$ and $\Delta_pu(x_0)=0$.
Consequently $\mathcal R_E(u)(x_0)=E|u(x_0)|^{p-2}u(x_0)$, and an eigenfunction with
eigenvalue $E$ cannot be of class $C^2$ near a critical point at which it does not
vanish; in particular, the first eigenfunction is not $C^2$ near its maximum points.
\end{lemma}

\begin{proof}
For $p>2$ the map $\Phi(\xi)=|\xi|^{p-2}\xi$ is of class $C^1(\R^d;\R^d)$ with
$D\Phi(0)=0$. Hence $\Phi(\nabla u)$ is $C^1$ and
$\Delta_pu(x_0)=\operatorname{tr}\big(D\Phi(\nabla u(x_0))D^2u(x_0)\big)=0$. If $u$ were a
$C^2$ eigenfunction near $x_0$, the equation would give
$0=-\Delta_pu(x_0)=E|u(x_0)|^{p-2}u(x_0)$ with $E>0$.
\end{proof}

\begin{remark}[Sources of error in the $p$-Laplacian experiments]\label{rem.plap.errors}
In addition to Remark \ref{rem.plap.lost}(3), the following facts explain why the
residuals and the errors are larger than for $p=2$.

(1) \emph{Limited regularity of the target.} By Lemma \ref{lem.critical} the
eigenfunctions are not $C^2$ at their nonzero critical points when $p>2$. For the radial
first eigenfunction in a ball one finds from \eqref{radial.ode} below
\[
u(r)=u(0)-\frac{p-1}{p}\Big(\frac{E\,u(0)^{p-1}}{d}\Big)^{\frac1{p-1}}r^{\frac p{p-1}}+o\big(r^{\frac p{p-1}}\big),
\]
so $u\in C^{1,1/(p-1)}$ and $|D^2u|\sim r^{-(p-2)/(p-1)}$ is unbounded at the origin if
$p>2$. The networks used here ($\tanh$ activation times a polynomial $B$) are smooth, so
at each of their critical points the pointwise residual equals $E|u|^{p-2}u$, which at
the maximum of the first eigenfunction is of the order of $E\|u\|_\infty^{p-1}$ and
cannot be reduced by training. The $L^2$ norm of the residual can still be made small,
but only by confining this defect to a small ball of radius $\varepsilon$ around the
critical point, which requires second derivatives of size
$\varepsilon^{-(p-2)/(p-1)}$ there; the contribution of the defect to the loss is of
order $E^2\|u\|_\infty^{2(p-1)}\varepsilon^{d}$. This difficulty is mild for $p$ close to
$2$ and becomes more severe as $p$ increases, and it is absent for $p=2$. Higher
eigenfunctions have more critical points, hence more such defects.

(2) \emph{Integrability of the strong residual.} For a smooth function $u$,
\[
\Delta_pu=|\nabla u|^{p-2}\Big(\Delta u+(p-2)\frac{\langle D^2u\nabla u,\nabla u\rangle}{|\nabla u|^2}\Big)
\qquad\text{where }\nabla u\ne0 .
\]
If $p\ge2$ this is bounded, so $\Lambda_p(u,E)<\infty$ for every network. If $1<p<2$,
$|\Delta_pu|$ behaves in general like $|x-x_0|^{p-2}$ near a nondegenerate critical point
$x_0$;
then $(\Delta_pu)^2$ is integrable near $x_0$ if and only if $p>2-\frac d2$ (always for
$d\ge2$, only for $p>\frac32$ if $d=1$), but its Monte Carlo estimator has finite variance
only if $p>2-\frac d4$ (that is, $p>\frac32$ for $d=2$). If $\nabla u$ vanishes
transversally on a hypersurface, $(\Delta_pu)^2$ is integrable only for $p>\frac32$.
Our experiments are restricted to $p>2$.

(3) \emph{Differentiation of $|\nabla u|^{p-2}\nabla u$.} The field
$\Phi(\nabla u_\theta)$ is differentiated automatically, once with respect to $x$ and
once more with respect to the parameters. The second operation involves
$D^2\Phi(\xi)=O(|\xi|^{p-3})$, which is unbounded at the critical points of $u_\theta$
when $p<3$: the gradient of the sampled loss then has heavy tails (its second moment is
finite only if $p>3-\frac d2$, i.e.\ $p>2$ for $d=2$), which penalizes the stochastic
optimization when $2<p<3$. A regularization
$(|\nabla u|^2+\varsigma^2)^{(p-2)/2}$ with $\varsigma>0$ fixed removes the singularity at
the price of a consistency error; in our implementation only a safeguard of size
$10^{-12}$ is added to $|\nabla u|$, so no consistency error is introduced, and the
singular set is hit with probability zero.

(4) \emph{No a posteriori control.} In the linear case the eigenvalue error is bounded by
$\eta^2/\gamma$; here no such bound exists and the error of the reported eigenvalue is of
the same order as the error of the eigenfunction, cf.\ Remark \ref{rem.plap.lost}(3).
The behaviour of the errors under refinement of the $E$-grid and under an increase of the
network size and of $N$ is reported in \texttt{[PLAP-TBD]}.
\end{remark}

\begin{remark}[A benchmark in any dimension: the first eigenvalue of a ball]\label{rem.radial}
The first eigenfunction of the ball $B_R\subset\R^d$ is radial (by simplicity of $E_1$
and rotation invariance), positive and decreasing, and solves
\begin{equation}\label{radial.ode}
\big(r^{d-1}|u'|^{p-2}u'\big)'+E\,r^{d-1}|u|^{p-2}u=0\quad\text{in }(0,R),\qquad u'(0)=0,\quad u(R)=0 .
\end{equation}
By homogeneity one may take $u(0)=1$, and $E_1(B_R)$ is the value of $E$ for which the
first zero of the solution of the initial value problem is $r=R$; it can be computed to
high accuracy by a shooting method applied to the first-order system
\[
u'=\operatorname{sign}(w)\,\big|w\,r^{1-d}\big|^{\frac1{p-1}},\qquad
w'=-E\,r^{d-1}|u|^{p-2}u,\qquad w=r^{d-1}|u'|^{p-2}u',
\]
started at $r=r_0\ll1$ with $w(r_0)=-Er_0^d/d$. Moreover, if $u$ is an eigenfunction in
$B_1$ with eigenvalue $E$, then $u(x/R)$ is an eigenfunction in $B_R$ with eigenvalue
$E/R^p$, so
\[
E_1(B_R)=\frac{E_1(B_1)}{R^p}.
\]
This provides reference values in dimensions where no numerical results are available
in the literature. We obtain
\[
\begin{array}{c|cccc}
(p,d) & (2.2,\,2) & (2.5,\,2) & (2.2,\,3) & (2.5,\,3)\\\hline
E_1(B_1) & 6.5317 & 7.7102 & 11.4699 & 14.1112
\end{array}
\]
(for $p=2$ the same code returns $5.78319$ and $9.86960$ for $d=2,3$; the value for
$(p,d)=(2.2,2)$ agrees with the value $6.5320$ of \cite{HJ11}). The higher zeros of the
shooting function give the radial eigenvalues, e.g.\ $41.7886$ for $(p,d)=(2.2,2)$; they
are eigenvalues but not the lowest ones: the second eigenfunction of the disk is not
radial \cite{Benedikt-Drabek-Girg}.
\end{remark}


%% ============================================================================
%% FRAGMENT 7 -- corrected proof detail for the Appendix (Theorem teo.min)
%% ============================================================================
%% In the proof of Theorem \ref{teo.min}, with ab <= eps a^2 + b^2/(4 eps) and
%% eps = 1/(4E) one gets 2E \int u \Delta u >= -(1/2)\int(\Delta u)^2 - 2E^2, hence
%% the correct form of \eqref{cota.Lambda2} is
\begin{equation*}
\Lambda(u_j,E)\ \ge\ \frac12\int_\Omega(\Delta u_j)^2\,dx-E^2 ,
\end{equation*}
%% (not "+E^2-E"); for E=0 the estimate is trivial.  The conclusion (boundedness
%% of \Delta u_j in L^2) is unchanged.  Boundedness in H^2 then uses the elliptic
%% estimate \|u\|_{H^2} <= C\|\Delta u\|_2 on H^2\cap H^1_0, valid for \Omega convex
%% or C^{1,1}; this hypothesis should be added to the statement.  The restriction
%% E >= 0 is not needed.  Alternatively the theorem follows at once from
%% \eqref{expansion}: the minimizers are the eigenfunctions of the eigenvalue(s)
%% closest to E.


%% ============================================================================
%% NEW BIBLIOGRAPHY ENTRIES (paper's style)
%% ============================================================================
%
% \bibitem{Anane1987}
% A.~Anane,
% Simplicit\'e et isolation de la premi\`ere valeur propre du $p$-laplacien avec poids,
% C. R. Acad. Sci. Paris S\'er. I Math. \textbf{305}, 725 (1987).
%
% \bibitem{Anane-Tsouli}
% A.~Anane and N.~Tsouli,
% On the second eigenvalue of the $p$-Laplacian,
% in \textit{Nonlinear Partial Differential Equations (F\`es, 1994)},
% Pitman Res. Notes Math. Ser. \textbf{343} (Longman, Harlow, 1996), pp.~1--9.
%
% \bibitem{Benedikt-Drabek-Girg}
% J.~Benedikt, P.~Dr\'abek, and P.~Girg,
% The second eigenfunction of the $p$-Laplacian on the disk is not radial,
% Nonlinear Anal. \textbf{75}, 4422 (2012).
%
% \bibitem{Cuesta-deFigueiredo-Gossez}
% M.~Cuesta, D.~de Figueiredo, and J.-P.~Gossez,
% The beginning of the Fu\v{c}ik spectrum for the $p$-Laplacian,
% J. Differential Equations \textbf{159}, 212 (1999).
%
% \bibitem{DeRyck-Lanthaler-Mishra}
% T.~De Ryck, S.~Lanthaler, and S.~Mishra,
% On the approximation of functions by tanh neural networks,
% Neural Networks \textbf{143}, 732 (2021).
%
% \bibitem{DiBenedetto}
% E.~DiBenedetto,
% $C^{1+\alpha}$ local regularity of weak solutions of degenerate elliptic equations,
% Nonlinear Anal. \textbf{7}, 827 (1983).
%
% \bibitem{E-Yu}
% W.~E and B.~Yu,
% The Deep Ritz method: a deep learning-based numerical algorithm for solving variational problems,
% Commun. Math. Stat. \textbf{6}, 1 (2018).
%
% \bibitem{GarciaAzorero-Peral}
% J.~P.~Garc\'ia Azorero and I.~Peral Alonso,
% Existence and nonuniqueness for the $p$-Laplacian: nonlinear eigenvalues,
% Comm. Partial Differential Equations \textbf{12}, 1389 (1987).
%
% \bibitem{Hornik1991}
% K.~Hornik,
% Approximation capabilities of multilayer feedforward networks,
% Neural Networks \textbf{4}, 251 (1991).
%
% \bibitem{Kato1949}
% T.~Kato,
% On the upper and lower bounds of eigenvalues,
% J. Phys. Soc. Japan \textbf{4}, 334 (1949).
%
% \bibitem{Lindqvist1990}
% P.~Lindqvist,
% On the equation $\operatorname{div}(|\nabla u|^{p-2}\nabla u)+\lambda|u|^{p-2}u=0$,
% Proc. Amer. Math. Soc. \textbf{109}, 157 (1990).
%
% \bibitem{Pinkus1999}
% A.~Pinkus,
% Approximation theory of the MLP model in neural networks,
% Acta Numer. \textbf{8}, 143 (1999).
%
% \bibitem{Rahaman-etal}
% N.~Rahaman, A.~Baratin, D.~Arpit, F.~Draxler, M.~Lin, F.~Hamprecht, Y.~Bengio, and A.~Courville,
% On the spectral bias of neural networks,
% in \textit{Proceedings of the 36th International Conference on Machine Learning},
% Proc. Mach. Learn. Res. \textbf{97}, 5301 (2019).
%
% \bibitem{Shin-Darbon-Karniadakis}
% Y.~Shin, J.~Darbon, and G.~E.~Karniadakis,
% On the convergence of physics informed neural networks for linear second-order elliptic
% and parabolic type PDEs,
% Commun. Comput. Phys. \textbf{28}, 2042 (2020).
%
% \bibitem{Tolksdorf}
% P.~Tolksdorf,
% Regularity for a more general class of quasilinear elliptic equations,
% J. Differential Equations \textbf{51}, 126 (1984).
%
% (Kato1949 is to be cited next to "Kato--Temple" in Proposition prop.rayleigh(c),
%  together with the already present \cite{Parlett1998}; E-Yu is to be cited where the
%  Rayleigh/Deep Ritz approach is mentioned in the cost comparison.)
